
\exo

Dans chacun des cas suivants, déterminer si les vecteurs $\vecu$ et $\vecv$ sont colinéaires.

\vspace{-6pt}
\setlength{\columnsep}{1cm}
\setlength{\columnseprule}{0pt}
\begin{multicols}{2}
\begin{enumerate}
\item $\vecu\couple{2}{3}$\; et \; $\vecv\couple{4}{6}$.$\hphantom{\frac{b}{q}}$
\medskip
\item $\vecu\couple{-3}{5}$\; et \; $\vecv\couple{9}{-15}$.$\hphantom{\frac{b}{q}}$
\medskip
\item $\vecu\couple{-3}{21}$\; et \; $\vecv\couple{4}{-28}$.$\hphantom{\frac{b}{q}}$
\medskip
\item $\vecu\couple{3}{12}$\; et \; $\vecv\couple{12}{3}$.$\hphantom{\frac{b}{q}}$
\item $\vecu\couple{\frac{7}{4}}{\frac{2}{7}}$\; et \; $\vecv\couple{\frac{1}{2}}{4}$.$\hphantom{\frac{b}{q}}$
\medskip
\item $\vecu\couple{\frac{2}{3}}{2}$\; et \; $\vecv\couple{\frac{1}{5}}{\frac{3}{5}}$.$\hphantom{\frac{b}{q}}$
\medskip
\item $\vecu\couple{-10}{4}$\; et \; $\vecv\couple{15}{-6}$.$\hphantom{\frac{b}{q}}$
\medskip
\item $\vecu\couple{\frac{3}{2}}{2}$\; et \; $\vecv\couple{21}{28}$.$\hphantom{\frac{b}{q}}$
\end{enumerate}
\end{multicols}
\vspace{-6pt}

\exo

\vspace{-6pt}
\setlength{\columnsep}{1cm}
\setlength{\columnseprule}{0pt}
\begin{multicols}{2}
\raggedcolumns
On note $\couple{x_1}{y_1}$ et $\couple{x_2}{y_2}$ les coordonnées de deux vecteurs $\vecu$ et $\vecv$.

Compléter la fonction Python ci-contre, de sorte qu'elle indique \pythoninline{Oui} si les vecteurs $\vecu$ et $\vecv$ sont colinéaires, et \pythoninline{Non} dans le cas contraire.

\columnbreak
\begin{pythoncodenumligne}
def colineaires(x1, y1, x2, y2):
    determinant = ................
    if determinant == ...........:
        return "Oui"
    else:
        ..........................
\end{pythoncodenumligne}
\end{multicols}
\vspace{-12pt}

\exo

Dans chacun des cas suivants, déterminer le nombre $z$ de sorte que les vecteurs $\vecu$ et $\vecv$ soient colinéaires.

\vspace{-6pt}
\setlength{\columnsep}{1cm}
\setlength{\columnseprule}{0pt}
\begin{multicols}{3}
\begin{enumerate}
\item $\vecu\couple{-3}{5}$\; et \; $\vecv\couple{z}{2}$.$\hphantom{\frac{b}{q}}$
\item $\vecu\couple{5}{z}$\; et \; $\vecv\couple{2}{\frac{1}{3}}$.$\hphantom{\frac{b}{q}}$
\item $\vecu\couple{-3}{7}$\; et \; $\vecv\couple{4}{z}$.$\hphantom{\frac{b}{q}}$
\end{enumerate}
\end{multicols}
\vspace{-6pt}

\exo

On considère les points $A\couple{0}{4}$,\; $B\couple{2}{7}$,\; $C\couple{8}{17}$ et $D\couple{16}{29}$.
\begin{enumerate}
\item
\begin{enumerate}
\item Déterminer les coordonnées des vecteurs $\vect{AB}$ et $\vect{CD}$.
\item Montrer que les vecteurs $\vect{AB}$ et $\vect{CD}$ sont colinéaires.
\end{enumerate}
\item Les vecteurs $\vect{AB}$ et $\vect{AC}$ sont-ils colinéaires ?
\end{enumerate}

\exo

\vspace{-6pt}
\setlength{\columnsep}{1cm}
\setlength{\columnseprule}{0pt}
\begin{multicols}{2}
\raggedcolumns
On a schématisé les rebonds d'une boule de billard. La boule est initialement au point $A\couple{2}{2,5}$, puis elle rebondit en $B\couple{4}{4}$, puis en $C\couple{8}{1}$, avant d'arriver au point $D$ d'ordonnée~$0$. On admet que le vecteur $\vect{CD}$ est colinéaire au vecteur $\vect{AB}$.

Calculer l'abscisse du point $D$.
\columnbreak
\begin{center}
\def\xmin{-1}
\def\xmax{9}
\def\ymin{-1}
\def\ymax{5}
\def\xunite{0.6cm}
\def\yunite{0.6cm}
\psset{unit=\xunite, xunit=\xunite, yunit=\yunite, algebraic=true, dimen=middle, dotstyle=*, dotsize=3pt 0, linewidth=0.8pt, arrowsize=3pt 2, arrowinset=0.25, comma=true}
\begin{pspicture*}(\xmin,\ymin)(\xmax,\ymax)
\psgrid[xunit=\xunite, yunit=\yunite, subgriddiv=0, gridlabels=0, gridcolor=lightgray](0,0)(\xmin,\ymin)(\xmax,\ymax)
\psaxes[labels=none,labelFontSize=\scriptstyle, xAxis=true, yAxis=true, Dx=1, Dy=1, ticksize=-2pt 0, subticks=1]{->}(0,0)(\xmin,\ymin)(\xmax,\ymax)[{\footnotesize $x$},135][{\footnotesize $y$},-45]
\pspolygon[fillcolor=black,fillstyle=solid,opacity=0.1](0,0)(8,0)(8,4)(0,4)
\psdot(2,2.5)
\psline[arrows=->,linewidth=1.2pt](2,2.5)(4,4)
\psline[arrows=->,linewidth=1.2pt](4,4)(8,1)
\psline[arrows=->,linewidth=1.2pt](8,1)(6.66,0)
\begin{footnotesize}
\uput{4pt}[180](2,2.5){$A$}
\uput{4pt}[090](4,4){$B$}
\uput{4pt}[000](8,1){$C$}
\uput{4pt}[270](6.66,0){$D$}
\uput{3pt}[270](1,0){$1$}
\uput{3pt}[180](0,1){$1$}
\uput{4pt}[225](0,0){$0$}
\end{footnotesize}
\end{pspicture*}
\end{center}
\end{multicols}

\exo

Dans chacun des cas suivants, déterminer si les points $A$, $B$ et $C$ sont alignés.

\vspace{-6pt}
\setlength{\columnsep}{1cm}
\setlength{\columnseprule}{0pt}
\begin{multicols}{2}
\begin{enumerate}
\item $A\couple{2}{13}$, $B\couple{-2}{-7}$ et $C\couple{11}{58}$.
\item $A\couple{9}{20}$, $B\couple{2}{-1}$ et $C\couple{25}{71}$.
\end{enumerate}
\end{multicols}
\vspace{-6pt}

\exo

Dans chacun des cas suivants, déterminer si le point $C$ appartient à la droite $(AB)$.

\vspace{-6pt}
\setlength{\columnsep}{1cm}
\setlength{\columnseprule}{0pt}
\begin{multicols}{2}
\begin{enumerate}
\item $A\couple{5}{-3}$, $B\couple{-3}{3}$ et $C\couple{15}{-9}$.
\item $A\couple{0}{-7}$, $B\couple{1}{0}$ et $C\couple{2}{7}$.
\end{enumerate}
\end{multicols}
\vspace{-6pt}

\exo

Dans chacun des cas suivants, déterminer si les droites $(AB)$ et $(CD)$ sont parallèles.

\vspace{-6pt}
\setlength{\columnsep}{1cm}
\setlength{\columnseprule}{0pt}
\begin{multicols}{2}
\begin{enumerate}
\item $A\couple{1}{2}$, $B\couple{5}{8}$, $C\couple{0}{-1}$ et $D\couple{5}{6}$.
\item $A\couple{3}{-10}$, $B\couple{15}{5}$, $C\couple{1}{1}$ et $D\couple{17}{21}$.
\end{enumerate}
\end{multicols}
\vspace{-6pt}

\exo

On note $y$ un réel, et on considère les points :
\[
A\couple{5}{-1}\qpvq B\couple{-1}{4}\qpvq C\couple{7}{2}\qpvq D\couple{1}{y}.
\]
\begin{enumerate}
\item Déterminer la valeur de $y$ de sorte que le point $D$ appartienne à la droite $(BC)$.
\item Déterminer la valeur de $y$ de sorte que les droites $(AB)$ et $(CD)$ soient parallèles.
\end{enumerate}

\exo

On considère les points :
\[
A\couple{4}{0}\qpvq B\couple{0}{7}\qpvq C\couple{-6}{-5}.
\]
\begin{enumerate}
\item Calculer les coordonnées du point $P$, milieu du segment $[AB]$.
\item Calculer les coordonnées des points $S$ et $T$ définis par :
\[
\vect{BS}=\tfrac{1}{3}\vect{BC}
\qetq
5\vect{CT}=4\vect{CA}.
\]
\item Le point $P$ appartient-il à la droite $(ST)$ ? Justifier par un calcul.
\end{enumerate}

\exo

\vspace{-6pt}
\setlength{\columnsep}{1cm}
\setlength{\columnseprule}{0pt}
\begin{multicols}{2}
\raggedcolumns
Dans un repère, on considère les points :
\begin{center}
$M\couple{7}{3}$, $N\couple{-3}{1}$, $C\couple{0}{5}$, $D\couple{5}{6}$ et $E\couple{3}{9}$.
\end{center}
\begin{enumerate}
\item Montrer que le quadrilatère $MNCD$ est un trapèze.
\item Montrer que $E$ est le point d'intersection des droites $(MD)$ et $(NC)$.
\item Calculer les coordonnées des points $J$ et $K$, milieux respectifs de $[MN]$ et $[CD]$.
\item Montrer que les points $E$, $J$ et $K$ sont alignés.
\end{enumerate}
\begin{center}
\def\xmin{-4.99}
\def\xmax{8}
\def\ymin{-0.99}
\def\ymax{10}
\def\xunite{0.5cm}
\def\yunite{0.5cm}
\psset{unit=\xunite, xunit=\xunite, yunit=\yunite, algebraic=true, dimen=middle, dotstyle=*, dotsize=3pt 0, linewidth=0.8pt, arrowsize=3pt 2, arrowinset=0.25, comma=true}
\begin{pspicture*}(\xmin,\ymin)(\xmax,\ymax)
%\psgrid[xunit=\xunite, yunit=\yunite, subgriddiv=0, gridlabels=0, gridcolor=lightgray](0,0)(\xmin,\ymin)(\xmax,\ymax)
\psaxes[labels=none,labelFontSize=\scriptstyle, xAxis=true, yAxis=true, Dx=1, Dy=1, ticksize=-2pt 0, subticks=1]{->}(0,0)(\xmin,\ymin)(\xmax,\ymax)[{\footnotesize $x$},135][{\footnotesize $y$},-45]
\psdots(7,3)(-3,1)(0,5)(5,6)(3,9)(2,2)(2.5,5.5)
\psline[linestyle=dashed, dash=2pt 2pt](7,3)(-3,1)
\psline[linestyle=dashed, dash=2pt 2pt](0,5)(5,6)
\psplot[linewidth=1.2pt, plotpoints=200]{\xmin}{\xmax}{1.33*x+5}
\psplot[linewidth=1.2pt, plotpoints=200]{\xmin}{\xmax}{-1.5*(x-3)+9}
\begin{footnotesize}
\uput{4pt}[040](7,3){$M$}
\uput{4pt}[040](5,6){$D$}
\uput{4pt}[150](0,5){$C$}
\uput{4pt}[150](-3,1){$N$}
\uput{4pt}[280](2,2){$J$}
\uput{4pt}[280](2.5,5.5){$K$}
\uput{4pt}[180](3,9){$E$}
\uput{4pt}[270](1,0){$1$}
\uput{4pt}[180](0,1){$1$}
\uput{5pt}[225](0,0){$0$}
\end{footnotesize}
\end{pspicture*}
\end{center}
\end{multicols}
\vspace{-6pt}

\exo

\vspace{-6pt}
\setlength{\columnsep}{1cm}
\setlength{\columnseprule}{0pt}
\begin{multicols}{2}
\raggedcolumns
Dans un repère, on considère les points :
\begin{center}
$A\couple{3}{4}$, $B\couple{1}{1}$, $C\couple{6}{-2}$ et $D\couple{8}{1}$.
\end{center}
Les points $I$, $E$ et $F$ sont tels que :
\[\vect{BI}=\tfrac{1}{2}\vect{BC},~\vect{AE}=\tfrac{1}{3}\vect{AC}\qetq \vect{CF}=\tfrac{1}{3}\vect{CA}.\]
\begin{enumerate}
\item Calculer les coordonnées des points $I$, $E$ et $F$.
\item 
\begin{enumerate}
\item Les vecteurs $\vect{BE}$ et $\vect{IF}$ sont-ils colinéaires ?
\item Que peut-on en déduire concernant les droites $(IE)$ et $(BF)$ ?
\end{enumerate}
\item Montrer que le quadrilatère $ABCD$ est un parallélogramme.
\item
\begin{enumerate}
\item Calculer la distance $AC$.
\item $ABCD$ est-il un rectangle ?
\end{enumerate}
\item Les points $I$, $F$ et $D$ sont-ils alignés ?
\end{enumerate}
\columnbreak
\begin{center}
\def\xmin{-1.99}
\def\xmax{9}
\def\ymin{-2.99}
\def\ymax{6}
\def\xunite{0.59cm}
\def\yunite{0.59cm}
\psset{unit=\xunite, xunit=\xunite, yunit=\yunite, algebraic=true, dimen=middle, dotstyle=*, dotsize=3pt 0, linewidth=0.8pt, arrowsize=3pt 2, arrowinset=0.25, comma=true}
\begin{pspicture*}(\xmin,\ymin)(\xmax,\ymax)
%\psgrid[xunit=\xunite, yunit=\yunite, subgriddiv=0, gridlabels=0, gridcolor=lightgray](0,0)(\xmin,\ymin)(\xmax,\ymax)
\psaxes[labels=none,labelFontSize=\scriptstyle, xAxis=true, yAxis=true, Dx=1, Dy=1, ticksize=-2pt 0, subticks=1]{->}(0,0)(\xmin,\ymin)(\xmax,\ymax)[{\footnotesize $x$},135][{\footnotesize $y$},-45]
\pspolygon[fillstyle=solid, fillcolor=lightgray,opacity=0.1](3,4)(1,1)(6,-2)
\psdots(3,4)(1,1)(6,-2)(8,1)(3.5,-0.5)(5,0)(4,2)
\begin{footnotesize}
\uput{4pt}[090](3,4){$A$}
\uput{4pt}[180](1,1){$B$}
\uput{4pt}[270](6,-2){$C$}
\uput{4pt}[090](8,1){$D$}
\uput{4pt}[050](4,2){$E$}
\uput{4pt}[050](5,0){$F$}
\uput{4pt}[240](3.5,-0.5){$I$}
\uput{4pt}[270](1,0){$1$}
\uput{4pt}[180](0,1){$1$}
\uput{5pt}[225](0,0){$0$}
\end{footnotesize}
\end{pspicture*}
\end{center}
\end{multicols}
